% =====================================================================
% frontmatter/cover.tex
% PT: Capa (página 1 do modelo ESTSetúbal/UNIPS). Reproduz a primeira
%     página de TrabalhoUNIPS_TecnologiaSetubal_Mestrado_definitiva.docx.
%     NÃO precisa de editar este ficheiro: o texto vem de
%     config/metadata.tex. NÃO altere as posições \put nem as medidas em
%     bp, porque copiam o modelo oficial do Word.
% EN: Cover (page 1 of the ESTSetúbal/UNIPS model). Reproduces the first
%     page of TrabalhoUNIPS_TecnologiaSetubal_Mestrado_definitiva.docx.
%     You do NOT need to edit this file: the text comes from
%     config/metadata.tex. Do NOT change the \put positions or the bp
%     measurements, because they copy the official Word model.
%
% PT: Como funciona
%     - Página A4 = 595 x 842 bp (1 bp = 1 ponto do Word = 1/72 pol.).
%     - Cada bloco é colocado com \put(x,-y) em relação ao canto superior
%       esquerdo da página: x para a direita, y para baixo (daí o sinal -).
%     - Fonte Arial (Helvetica, com as mesmas métricas), sem hifenização.
%     - No Word, cada caixa de texto tem margem interna de 7.2 bp à
%       esquerda e 3.6 bp no topo; por isso cada \put fica em
%       x = caixa + 7.2 e y = topo da caixa + 3.6 + ascendente da Arial,
%       ou seja, na linha de base da primeira linha do bloco.
% EN: How it works
%     - A4 page = 595 x 842 bp (1 bp = 1 Word point = 1/72 in).
%     - Each block is placed with \put(x,-y) relative to the top-left
%       corner of the page: x to the right, y downwards (hence the - sign).
%     - Arial font (Helvetica, with the same metrics), no hyphenation.
%     - In Word each text box has an inner margin of 7.2 bp on the left and
%       3.6 bp at the top; so each \put sits at x = box + 7.2 and
%       y = box top + 3.6 + Arial ascender, i.e. on the baseline of the
%       first line of the block.
%
% PT: Medidas das caixas do Word (bp, a partir do canto superior esquerdo)
% EN: Word text-box measurements (bp, from the top-left corner)
%   logotipo / logo        x = 191.05, topo/top y = 67, largura/width 255.4
%   autor / author         x = 85.15,  y = 263.4, largura/width 160.0, 16 pt, à direita/right-aligned
%   título / title         x = 249.95, y = 263.0, largura/width 290.1, 16 pt negrito/bold
%   grau / degree          x = 249.95, y = 460.0, largura/width 275.1, 12 pt
%   linha / line           centro/centre x = 250.65, y = 263.5 -> 773.7, 3.75 pt, #CAF4FF
%   orientadores / supervisors  x = 250.25, y = 582.2, largura/width 275.1, 12 pt
%   data / date            x = 250.25, y = 735.5, 10 pt
%
% PT: Título longo
%     O título (16 pt, entrelinha 18.4 bp) começa na linha de base
%     y = 278.6 e o bloco do grau começa em y = 472.0. Cabem no máximo
%     10 linhas de título (cerca de 20 a 28 caracteres em maiúsculas por
%     linha, ou seja, uns 220 caracteres no total); a 11.ª linha sobrepõe-se
%     ao bloco do grau. Para um aspeto equilibrado, mantenha o título até
%     8 linhas. O subtítulo conta para estas linhas. Se não couber,
%     encurte o título; pode também forçar as quebras de linha com \\ no
%     próprio \thesistitle. A folha do júri aceita até 12 linhas, por isso
%     o limite é dado pela capa.
% EN: Long title
%     The title (16 pt, 18.4 bp leading) starts on baseline y = 278.6 and
%     the degree block starts at y = 472.0. At most 10 title lines fit
%     (roughly 20 to 28 capital letters per line, about 220 characters in
%     total); an 11th line overlaps the degree block. For a balanced look,
%     keep the title within 8 lines. The subtitle counts towards these
%     lines. If it does not fit, shorten the title; you can also force
%     line breaks with \\ inside \thesistitle. The jury page accepts up to
%     12 lines, so the cover sets the limit.
% =====================================================================
\begin{titlepage}
% PT: Os blocos são desenhados por cima da página, no momento em que ela é
%     enviada para o PDF, com coordenadas absolutas.
% EN: The blocks are drawn on top of the page, when it is shipped out to
%     the PDF, with absolute coordinates.
\AddToHookNext{shipout/foreground}{%
  \setlength{\unitlength}{1bp}%
  \linespread{1}\fontfamily{phv}\fontseries{m}\fontshape{n}\selectfont
  \hyphenpenalty=10000 \exhyphenpenalty=10000
  % -------------------------------------------------------------------
  % PT: Logotipo da ESTSetúbal/UNIPS (figures/estsetubal_logo.pdf), centrado no topo.
  % EN: ESTSetúbal/UNIPS logo (figures/estsetubal_logo.pdf), centred at the top.
  \put(191.05,-67){\makebox(0,0)[lt]{\includegraphics[width=255.4bp]{figures/estsetubal_logo.pdf}}}%
  % -------------------------------------------------------------------
  % PT: Autor: coluna da esquerda, 16 pt regular, maiúsculas, alinhado à direita.
  % EN: Author: left column, 16 pt regular, capitals, right-aligned.
  \put(92.35,-279.0){\parbox[t]{145.6bp}{\raggedleft\fontsize{16bp}{18.4bp}\selectfont
    \MakeUppercase{\authorname}}}%
  % -------------------------------------------------------------------
  % PT: Título (e subtítulo opcional): coluna da direita, 16 pt negrito, maiúsculas.
  % EN: Title (and optional subtitle): right column, 16 pt bold, capitals.
  \put(257.15,-278.6){\parbox[t]{275.7bp}{\raggedright\fontsize{16bp}{18.4bp}\selectfont\bfseries
    \MakeUppercase{\thesistitle}\ifdefempty{\thesissubtitle}{}{\MakeUppercase{: \thesissubtitle}}}}%
  % -------------------------------------------------------------------
  % PT: Grau: duas linhas, 12 pt regular (\coverdegree e \coursefield).
  % EN: Degree: two lines, 12 pt regular (\coverdegree and \coursefield).
  \put(257.15,-472.0){\parbox[t]{260.7bp}{\raggedright\fontsize{12bp}{13.8bp}\selectfont
    \coverdegree\par\coursefield}}%
  % -------------------------------------------------------------------
  % PT: Linha azul vertical #CAF4FF entre as duas colunas (no Word é uma
  %     linha de 510.2 pt rodada 270 graus).
  % EN: Vertical blue line #CAF4FF between the two columns (in Word it is a
  %     510.2 pt line rotated 270 degrees).
  \put(248.775,-773.7){\textcolor[HTML]{CAF4FF}{\rule{3.75bp}{510.2bp}}}%
  % -------------------------------------------------------------------
  % PT: Orientador e coorientador, 12 pt. O bloco COORIENTADOR só aparece
  %     se \cosupervisorname não estiver vazio.
  % EN: Supervisor and co-supervisor, 12 pt. The COORIENTADOR block only
  %     appears if \cosupervisorname is not empty.
  \put(257.45,-594.2){\parbox[t]{260.7bp}{\raggedright\fontsize{12bp}{13.8bp}\selectfont
    \setlength{\parskip}{2bp}\setlength{\parindent}{0pt}%
    {\bfseries \supervisorlabel\par}
    \supervisorname\par
    \ifdefempty{\cosupervisorname}{}{%
      \vspace{6bp}
      {\bfseries \cosupervisorlabel\par}
      \cosupervisorname\par}}}%
  % -------------------------------------------------------------------
  % PT: Data da realização da prova (mês e ano), 10 pt.
  % EN: Date of the public defence (month and year), 10 pt.
  \put(257.45,-747.4){\parbox[t]{230.5bp}{\raggedright\fontsize{10bp}{13.2bp}\selectfont
    \defencemonth\ de \defenceyear}}%
}
% PT: Página vazia: o conteúdo é todo desenhado pelo bloco acima.
% EN: Empty page: all content is drawn by the block above.
\null
\end{titlepage}
